% Plain-English seminar slides for a general audience.
% Same theme and figure path as uk_long_rate_seminar.tex.
% The full deck is left unchanged.
%
% Compile from this directory:
%   pdflatex uk_long_rate_simple.tex
%   latexmk -pdf uk_long_rate_simple.tex
\documentclass[11pt,aspectratio=169]{beamer}

\usetheme{Madrid}
\usecolortheme{default}
\setbeamertemplate{navigation symbols}{}
\setbeamertemplate{itemize items}[circle]
\setbeamerfont{footnote}{size=\tiny}

\usepackage[T1]{fontenc}
\usepackage[utf8]{inputenc}
\usepackage{amsmath,amssymb}
\usepackage{booktabs}
\usepackage{graphicx}
\usepackage{array}

\newcommand{\irfdir}{../output/slides/}
\graphicspath{{\irfdir}}

\title[Long-term UK borrowing costs]{When long-term UK borrowing costs rise}
\subtitle{And Bank Rate stays put}
\author[UK long-rate]{A short guide to the UK long-rate model}
\institute[MacroModellerBot]{\texttt{uk-long-rate/} in MacroModellerBot}
\date{September 2026}

\begin{document}

\begin{frame}[plain]
  \titlepage
\end{frame}

\begin{frame}{The question}
  What happens when long-term UK borrowing costs rise,
  while Bank Rate stays put?
  \vspace{0.6em}
  \begin{itemize}
    \item Long-term rates are what mortgages, firms, and the government
          pay to borrow for years, not the overnight policy rate.
    \item They can rise because investors demand extra compensation
          for holding long debt, even if the Bank does not change
          Bank Rate.
    \item This talk looks at that case: a one-off rise of 100 basis
          points in long rates, with Bank Rate held on its baseline
          path for three years.
    \item Every result below is that experiment, measured against a
          baseline in which the rise does not happen.
  \end{itemize}
\end{frame}

\begin{frame}{How it spreads}
  \textbf{In the model}
  \begin{itemize}
    \item \textbf{Mortgages.} Most deals are fixed. The higher rate
          bites as those deals roll off, not all at once.
    \item \textbf{House prices.} Dearer new mortgages pull house prices
          down, which weakens borrowers who rely on the house as
          collateral.
    \item \textbf{Business investment.} A higher long-term cost of
          capital leads firms to invest less.
    \item \textbf{The government interest bill.} New and refinanced
          gilts cost more. Part of the stock is index-linked, so the
          bill does not jump immediately.
    \item \textbf{Sterling.} With Bank Rate held fixed, the pound
          strengthens in this experiment.
  \end{itemize}
  \vspace{0.35em}
  \textbf{Described, but not modelled}
  \begin{itemize}
    \item Pension-fund (LDI) stress, as in September 2022.
    \item A sovereign-risk shock, in which gilt worries also hit
          the pound.
  \end{itemize}
\end{frame}

\begin{frame}{Why timing matters}
  \begin{itemize}
    \item Only about 9\% of the mortgage stock reprices each quarter.
          That share is $\lambda_q = 0.09$.
    \item A borrower whose fix has years left does not pay the new
          rate yet. Someone remortgaging this quarter does.
    \item So the average mortgage rate on the whole stock climbs
          slowly, even when the rate on new loans jumps.
    \item The pressure on household spending therefore builds over
          several quarters, rather than arriving in one month.
  \end{itemize}
\end{frame}

\begin{frame}{What the model says}
  \begin{center}
    \includegraphics[width=\textwidth,height=0.42\textheight,keepaspectratio]%
      {\irfdir a_macro.pdf}
  \end{center}
  \vspace{-0.15em}
  {\scriptsize
  Response to a one-off $+100$~bp rise in long rates, measured against
  a baseline without it.}
  \begin{itemize}
    \item Output is 0.55\% lower than it would otherwise be
          (trough at quarter 4).
    \item House prices are 2.6\% lower than otherwise (quarter 8).
    \item The government interest bill is 0.20~pp of GDP higher than
          otherwise (peak at quarter 24).
  \end{itemize}
\end{frame}

\begin{frame}{Who is hit}
  \begin{center}
    \includegraphics[width=\textwidth,height=0.34\textheight,keepaspectratio]%
      {\irfdir a_consumption_mortgage.pdf}
  \end{center}
  \vspace{-0.2em}
  {\scriptsize
  Same one-off $+100$~bp rise in long rates, against a baseline
  without it.}
  \begin{itemize}
    \item Savers cut back first: 0.77\% below baseline at quarter 3.
    \item Borrowers cut back more, and later: 1.04\% below baseline
          at quarter 9, as their mortgages reprice.
    \item That lag is the 9\% quarterly reset. The rate on new loans
          jumps immediately; the rate on the whole stock does not.
  \end{itemize}
\end{frame}

\begin{frame}{If markets instead expect Bank Rate to rise}
  \begin{center}
    \includegraphics[width=\textwidth,height=0.32\textheight,keepaspectratio]%
      {\irfdir a_vs_b.pdf}
  \end{center}
  \vspace{-0.15em}
  {\scriptsize
  Both lines are a one-off $+100$~bp rise in long rates, against a
  baseline without it. Only the solid line keeps Bank Rate on hold.}
  \begin{itemize}
    \item If the long rate rises because Bank Rate itself is expected
          to rise, Bank Rate is about 315~bp higher than otherwise
          on impact, and about 560~bp higher at quarter 4.
    \item Output is then 1.178\% lower than otherwise (quarter 5),
          and house prices 6.551\% lower (quarter 7).
    \item That is a deeper hit than when Bank Rate stays put
          (0.55\% and 2.6\%).
  \end{itemize}
\end{frame}

\begin{frame}{Caveats}
  \begin{itemize}
    \item The numbers come from a calibrated model. They are not
          estimates fitted to UK data.
    \item A few settings are provisional choices, including how
          strongly inflation responds and how house prices are
          determined.
    \item The sovereign-risk scenario is not yet modelled.
    \item LDI stress is not modelled. The September 2022 gilt spiral
          is a separate, non-linear event.
  \end{itemize}
\end{frame}

\begin{frame}{Next steps}
  \begin{itemize}
    \item Estimate the effects directly from high-frequency surprises
          in gilt yields, using local projections.
    \item Include a state-dependent version, so the effect can differ
          when the gilt market is stressed.
    \item Then estimate this model on UK data, with a wide prior on
          how strongly inflation responds.
    \item The sovereign-risk scenario is still to be built.
  \end{itemize}
\end{frame}

\begin{frame}{Appendix: three fixed settings}
  \begin{itemize}
    \item About 9\% of mortgages reprice each quarter
          ($\lambda_q = 0.09$).
    \item A quarter of the gilt stock is index-linked (share 0.25).
    \item Average gilt duration is 9.1 years, so the government's
          coupon book refixes slowly.
  \end{itemize}
  \vspace{0.4em}
  The rate on the whole mortgage stock is mostly last quarter's rate:
  \[
    R^{\mathrm{stock}}_{t}
      = (1-\lambda_q)\, R^{\mathrm{stock}}_{t-1}
      + \lambda_q\, R^{\mathrm{new}}_{t}.
  \]
  {\scriptsize With $\lambda_q = 0.09$, nine tenths of the stock keeps
  the old rate each quarter.}
\end{frame}

\end{document}
